% =====================================================================
%  Chapter 2 — Background and Literature
%  Input with:  % =====================================================================
%  Chapter 2 — Background and Literature
%  Input with:  % =====================================================================
%  Chapter 2 — Background and Literature
%  Input with:  % =====================================================================
%  Chapter 2 — Background and Literature
%  Input with:  \input{chapters/02_background}
%  No tables or figures.
% =====================================================================

\section{Background and Literature}
\label{sec:background}

\subsection{A common institutional inheritance, and an uncommon starting point}
\label{sec:soviet-inheritance}

Until December 1991 the fourteen countries studied here were administrative units of a single state. They shared a health system, a school curriculum, a language of government, a pension architecture and a family-policy regime, including maternity leave, child benefits and the near-free availability of abortion that made the Soviet Union unusual among industrial societies. If institutional convergence produced demographic convergence, these are the countries where one would expect to see it.

They did not start from a common position, though, and the reason matters for what follows. Fertility in Central Asia through the 1920s and 1930s fluctuated around four to five children per woman, on a par with much of the pre-transitional world. What happened next runs against the usual direction of demographic change. From the early 1940s fertility in the region \emph{rose}, continuing to climb until the early 1960s and holding on a plateau until the early 1970s, before beginning a long decline that bottomed out around 2000 \parencite{spoorenberg2017}.

This happened during exactly the decades when Soviet policy was extending female education, drawing women into the labour force and urbanising the population, the conditions under which the standard account expects fertility to fall. \textcite{spoorenberg2017} rules out the most obvious explanation, showing that marriage was being postponed rather than accelerated during the increase: in Kazakhstan the share of women married at ages 20--24 fell from 78 per cent in 1939 to 62 per cent in 1959. What appears to have driven the rise instead is the removal of biological and behavioural checks on childbearing. Life expectancy in the region rose from around 30--35 years in the mid-1920s to 50--55 by the early 1950s; sterility and infant loss fell with the extension of the health system; and urban migration shortened breastfeeding and therefore birth intervals. Development raised fertility by unleashing childbearing capacity before it began to suppress it, a pattern documented in other early-transition settings \parencite{nag1980,romaniuk1980}.

Two things follow. Central Asia's high fertility at the end of the Soviet period was not simply a pre-modern residue waiting to be modernised away; it was itself partly a product of twentieth-century development. And the region entered independence at a very different point in its fertility history from the European republics, which had been below or near replacement for decades. Whatever else 1991 did, it did not apply a common shock to a common starting position.

\subsection{The transition shock and how it has been explained}
\label{sec:transition}

Fertility fell everywhere in the region after 1991, and in most places it fell hard. Between 1991 and 2000 the total fertility rate declined by 38.6 per cent in Uzbekistan, 33.3 per cent in Kyrgyzstan and 32 per cent in Kazakhstan \parencite{spoorenberg2013}; the European republics fell further and from lower levels, with fifteen post-communist countries recording a rate below 1.3 at least once by 2004 \parencite{billingsley2010}. Populations shrank accordingly. Between 1989 and 2008 Georgia lost nearly 19 per cent of its population and Moldova nearly 18, while Latvia and Estonia lost around 15 per cent each. Over the same period Tajikistan grew by 42 per cent and Uzbekistan by 34 \parencite{brainerd2010}. The demographic divergence this thesis measures in fertility rates was already visible in population totals a decade before the study window opens.

\paragraph{The economic account.} The standard economic framework treats children as time-intensive, so that the shadow price of a child rises with the mother's wage \parencite{becker1981}. On this reading the post-communist decline reflects a rising opportunity cost of childbearing: women's relative wages improved in much of Eastern Europe, and returns to education rose across nearly all transition countries, drawing women into tertiary enrolment and delaying family formation \parencite{brainerd2010}. The account is incomplete on its own terms. \textcite{brainerd2010} notes that women's wages did not improve in Russia or Ukraine, where fertility nonetheless collapsed, and the framework says nothing about uncertainty, which is what the transition mostly supplied. \textcite{kohler2002} argue that the Russian decline of the early and mid-1990s was driven by labour-market crisis and the option value of postponing an irreversible commitment under uncertain conditions. Delay is cheap when the future is unreadable.

\paragraph{Three competing explanations.} \textcite{billingsley2010} organises the debate into three accounts that differ in economic context, in the fertility behaviour they predict, and in what motivates it. The second demographic transition attributes decline to ideational change and self-actualisation, and expects it under economic stability. The postponement transition attributes it to delay of first births as a rational response to uncertainty. The economic-crisis account predicts something different in kind: not delay but stopping, the abandonment of higher-order births under material pressure.

The distinction matters because the two behaviours leave different traces. Postponement depresses the period fertility rate temporarily and may be recovered later; stopping reduces completed family size. Billingsley's central finding is that no single account fits the whole region: postponement explains the bulk of the decline in no more than five countries, all of them in Central Europe, while elsewhere much of the fall happened before postponement began. She also finds the two processes tied to economic conditions in opposite directions, improving conditions with postponement and deteriorating conditions with stopping.

One of her observations is worth carrying forward. Postponement in the post-communist region began from a much younger base than in Southern Europe: mothers' average age at first birth in 1999 was 23.9 years against 28.3 in the Southern European countries where the postponement literature was developed. Delay starting at 24 need not translate into births forgone, which is a caution against reading period declines in this region as permanent.

\paragraph{Policy.} Governments across the region responded with pronatalist measures, of which Russia's 2006 maternity-capital programme is the largest and most studied \parencite{brainerd2010}. \textcite{zakharov2024} assesses it over a long horizon and finds the effect narrower than its reputation: Russian fertility began rising in 2000, six years before the decree, and what the policy most visibly changed was the interval between births rather than the number of them. Period rates rose while births were brought forward, then fell back to 1.50 by 2019, essentially the pre-policy level. He reads this as the same inflation-and-correction pattern produced by the pronatalist campaign of the 1980s, and concludes that two decades of policy did not halt the underlying modernisation of Russian fertility. The wider evidence points the same way: financial incentives tend to move the timing of childbearing more than its quantum \parencite{gauthier2007}.

\subsection{What is known about Central Asia}
\label{sec:ca-literature}

The Central Asian literature is smaller, more recent, and built almost entirely on country cases rather than regional comparisons.

The starting point is that the fertility recovery after 2000 is real. \textcite{spoorenberg2015} tests the possibility that it reflects improving completeness of birth registration by cross-comparing vital-registration rates against census and survey estimates computed by independent methods, and rules the artefact explanation out, with a reservation for Turkmenistan, where independent data are not accessible. This matters more than it might appear: in a region where one country has not held a census since 1989, establishing that a demographic movement is not a recording artefact is a precondition for studying it at all.

Beyond that, four strands of explanation have developed, and because Chapter~\ref{sec:discussion} works through them in detail, they are only mapped here.

The first is population composition. Fertility differs by about one child per woman between women of the titular ethnicities and women of European origin, consistently across countries and data sources, and post-1991 emigration changed the ethnic composition of these populations substantially \parencite{spoorenberg2013,agadjanian2013}. Holding ethnic-specific fertility constant and letting composition alone change reproduces a large share of the observed national increases \parencite{spoorenberg2015}. The second is marriage: near-universal but with rising age at first union, and sharply differentiated by education \parencite{dommaraju2008}. The third is culture, understood not as individual belief but as group-level norms and expectations. \textcite{agadjanian2022} distinguish the normative propensities a group transmits from its position in a society's hierarchy, and find the two predicting different things: completed fertility follows the first, desires for further children the second. \textcite{kumo2024} approach the same territory from the religion side, modelling a chain from religiosity through gender-role beliefs to family size and asking how far Soviet institutions interrupted it. The fourth is the economic context, which \textcite{spoorenberg2015} finds positively associated with fertility in a five-country panel, alongside a shifting tempo effect accounting for part of the rise in Kyrgyzstan and rather less of it in Kazakhstan. Labour migration belongs here too: remittances reach a fifth or more of GDP in Tajikistan and Kyrgyzstan, and their relationship with fertility across the post-communist world is negative on average but heterogeneous \parencite{tokhirov2024}.

Two features of this literature are worth naming. It is predominantly micro-level, working with survey and census microdata, which is what allows it to identify mechanisms a country-level design cannot. And it is largely internal to the region, comparing Central Asian countries with one another or ethnic groups within them rather than against the rest of the post-Soviet space. \textcite{nedoluzhko2012} is the most explicitly comparative, and finds country of residence mattering more for fertility than ethnicity, even among co-ethnics living either side of a Central Asian border.

\subsection{The gap}
\label{sec:gap-literature}

Three specific gaps motivate what follows.

First, the comparative post-communist literature has largely left Central Asia out of its formal analysis. \textcite{billingsley2010} includes all five republics in her descriptive treatment but restricts the regression sample, excluding the majority of the Central Asian republics; of the five, only Kazakhstan enters the statistical analysis \parencite{spoorenberg2015}. The result is that the region most likely to falsify a general account of post-communist fertility is the region least represented in the tests of it.

Second, the literature stops too early. Billingsley's window closes in 2003, Spoorenberg's principal analyses run to 2009--2011, and the ethnic-differential studies rest on surveys from the 2000s. The turning point documented in Chapter~\ref{sec:trends}, a reversal in dispersion around 2016 and a widening difference thereafter, falls entirely outside every one of these windows. The Uzbek fertility surge after 2017, the largest single movement in the panel assembled here, postdates almost all of the comparative work on the region; the one substantial study of it is a national survey commissioned precisely because existing theory did not anticipate it \parencite{unfpa2023}.

Third, the two literatures use different units and have not been joined. The Central Asian work is micro-level and country-specific; the post-communist comparative work is macro-level and mostly excludes Central Asia. Where a macro treatment of Central Asian fertility exists, it is confined to the region and uses a single economic covariate: \textcite{spoorenberg2015} relates period fertility to GDP per capita in a five-country panel and identifies the reliance on one variable as a limitation. No study places all fourteen comparable post-Soviet states in a single frame, over a window reaching to 2023, with a set of economic controls and inference appropriate to a sample of this size.

This thesis does that, and the ambition is deliberately bounded. It does not attempt to identify why Central Asian fertility is higher; Section~\ref{sec:interpretation} explains why a fourteen-country panel cannot. What it can establish is how large the difference is, whether it has been closing, how much of it moves together with the macroeconomic conditions the comparative literature relies on, and how much is left over for the mechanisms the micro literature describes. Setting that boundary precisely is the contribution, and the size of what remains outside it is the result.

%  No tables or figures.
% =====================================================================

\section{Background and Literature}
\label{sec:background}

\subsection{A common institutional inheritance, and an uncommon starting point}
\label{sec:soviet-inheritance}

Until December 1991 the fourteen countries studied here were administrative units of a single state. They shared a health system, a school curriculum, a language of government, a pension architecture and a family-policy regime, including maternity leave, child benefits and the near-free availability of abortion that made the Soviet Union unusual among industrial societies. If institutional convergence produced demographic convergence, these are the countries where one would expect to see it.

They did not start from a common position, though, and the reason matters for what follows. Fertility in Central Asia through the 1920s and 1930s fluctuated around four to five children per woman, on a par with much of the pre-transitional world. What happened next runs against the usual direction of demographic change. From the early 1940s fertility in the region \emph{rose}, continuing to climb until the early 1960s and holding on a plateau until the early 1970s, before beginning a long decline that bottomed out around 2000 \parencite{spoorenberg2017}.

This happened during exactly the decades when Soviet policy was extending female education, drawing women into the labour force and urbanising the population, the conditions under which the standard account expects fertility to fall. \textcite{spoorenberg2017} rules out the most obvious explanation, showing that marriage was being postponed rather than accelerated during the increase: in Kazakhstan the share of women married at ages 20--24 fell from 78 per cent in 1939 to 62 per cent in 1959. What appears to have driven the rise instead is the removal of biological and behavioural checks on childbearing. Life expectancy in the region rose from around 30--35 years in the mid-1920s to 50--55 by the early 1950s; sterility and infant loss fell with the extension of the health system; and urban migration shortened breastfeeding and therefore birth intervals. Development raised fertility by unleashing childbearing capacity before it began to suppress it, a pattern documented in other early-transition settings \parencite{nag1980,romaniuk1980}.

Two things follow. Central Asia's high fertility at the end of the Soviet period was not simply a pre-modern residue waiting to be modernised away; it was itself partly a product of twentieth-century development. And the region entered independence at a very different point in its fertility history from the European republics, which had been below or near replacement for decades. Whatever else 1991 did, it did not apply a common shock to a common starting position.

\subsection{The transition shock and how it has been explained}
\label{sec:transition}

Fertility fell everywhere in the region after 1991, and in most places it fell hard. Between 1991 and 2000 the total fertility rate declined by 38.6 per cent in Uzbekistan, 33.3 per cent in Kyrgyzstan and 32 per cent in Kazakhstan \parencite{spoorenberg2013}; the European republics fell further and from lower levels, with fifteen post-communist countries recording a rate below 1.3 at least once by 2004 \parencite{billingsley2010}. Populations shrank accordingly. Between 1989 and 2008 Georgia lost nearly 19 per cent of its population and Moldova nearly 18, while Latvia and Estonia lost around 15 per cent each. Over the same period Tajikistan grew by 42 per cent and Uzbekistan by 34 \parencite{brainerd2010}. The demographic divergence this thesis measures in fertility rates was already visible in population totals a decade before the study window opens.

\paragraph{The economic account.} The standard economic framework treats children as time-intensive, so that the shadow price of a child rises with the mother's wage \parencite{becker1981}. On this reading the post-communist decline reflects a rising opportunity cost of childbearing: women's relative wages improved in much of Eastern Europe, and returns to education rose across nearly all transition countries, drawing women into tertiary enrolment and delaying family formation \parencite{brainerd2010}. The account is incomplete on its own terms. \textcite{brainerd2010} notes that women's wages did not improve in Russia or Ukraine, where fertility nonetheless collapsed, and the framework says nothing about uncertainty, which is what the transition mostly supplied. \textcite{kohler2002} argue that the Russian decline of the early and mid-1990s was driven by labour-market crisis and the option value of postponing an irreversible commitment under uncertain conditions. Delay is cheap when the future is unreadable.

\paragraph{Three competing explanations.} \textcite{billingsley2010} organises the debate into three accounts that differ in economic context, in the fertility behaviour they predict, and in what motivates it. The second demographic transition attributes decline to ideational change and self-actualisation, and expects it under economic stability. The postponement transition attributes it to delay of first births as a rational response to uncertainty. The economic-crisis account predicts something different in kind: not delay but stopping, the abandonment of higher-order births under material pressure.

The distinction matters because the two behaviours leave different traces. Postponement depresses the period fertility rate temporarily and may be recovered later; stopping reduces completed family size. Billingsley's central finding is that no single account fits the whole region: postponement explains the bulk of the decline in no more than five countries, all of them in Central Europe, while elsewhere much of the fall happened before postponement began. She also finds the two processes tied to economic conditions in opposite directions, improving conditions with postponement and deteriorating conditions with stopping.

One of her observations is worth carrying forward. Postponement in the post-communist region began from a much younger base than in Southern Europe: mothers' average age at first birth in 1999 was 23.9 years against 28.3 in the Southern European countries where the postponement literature was developed. Delay starting at 24 need not translate into births forgone, which is a caution against reading period declines in this region as permanent.

\paragraph{Policy.} Governments across the region responded with pronatalist measures, of which Russia's 2006 maternity-capital programme is the largest and most studied \parencite{brainerd2010}. \textcite{zakharov2024} assesses it over a long horizon and finds the effect narrower than its reputation: Russian fertility began rising in 2000, six years before the decree, and what the policy most visibly changed was the interval between births rather than the number of them. Period rates rose while births were brought forward, then fell back to 1.50 by 2019, essentially the pre-policy level. He reads this as the same inflation-and-correction pattern produced by the pronatalist campaign of the 1980s, and concludes that two decades of policy did not halt the underlying modernisation of Russian fertility. The wider evidence points the same way: financial incentives tend to move the timing of childbearing more than its quantum \parencite{gauthier2007}.

\subsection{What is known about Central Asia}
\label{sec:ca-literature}

The Central Asian literature is smaller, more recent, and built almost entirely on country cases rather than regional comparisons.

The starting point is that the fertility recovery after 2000 is real. \textcite{spoorenberg2015} tests the possibility that it reflects improving completeness of birth registration by cross-comparing vital-registration rates against census and survey estimates computed by independent methods, and rules the artefact explanation out, with a reservation for Turkmenistan, where independent data are not accessible. This matters more than it might appear: in a region where one country has not held a census since 1989, establishing that a demographic movement is not a recording artefact is a precondition for studying it at all.

Beyond that, four strands of explanation have developed, and because Chapter~\ref{sec:discussion} works through them in detail, they are only mapped here.

The first is population composition. Fertility differs by about one child per woman between women of the titular ethnicities and women of European origin, consistently across countries and data sources, and post-1991 emigration changed the ethnic composition of these populations substantially \parencite{spoorenberg2013,agadjanian2013}. Holding ethnic-specific fertility constant and letting composition alone change reproduces a large share of the observed national increases \parencite{spoorenberg2015}. The second is marriage: near-universal but with rising age at first union, and sharply differentiated by education \parencite{dommaraju2008}. The third is culture, understood not as individual belief but as group-level norms and expectations. \textcite{agadjanian2022} distinguish the normative propensities a group transmits from its position in a society's hierarchy, and find the two predicting different things: completed fertility follows the first, desires for further children the second. \textcite{kumo2024} approach the same territory from the religion side, modelling a chain from religiosity through gender-role beliefs to family size and asking how far Soviet institutions interrupted it. The fourth is the economic context, which \textcite{spoorenberg2015} finds positively associated with fertility in a five-country panel, alongside a shifting tempo effect accounting for part of the rise in Kyrgyzstan and rather less of it in Kazakhstan. Labour migration belongs here too: remittances reach a fifth or more of GDP in Tajikistan and Kyrgyzstan, and their relationship with fertility across the post-communist world is negative on average but heterogeneous \parencite{tokhirov2024}.

Two features of this literature are worth naming. It is predominantly micro-level, working with survey and census microdata, which is what allows it to identify mechanisms a country-level design cannot. And it is largely internal to the region, comparing Central Asian countries with one another or ethnic groups within them rather than against the rest of the post-Soviet space. \textcite{nedoluzhko2012} is the most explicitly comparative, and finds country of residence mattering more for fertility than ethnicity, even among co-ethnics living either side of a Central Asian border.

\subsection{The gap}
\label{sec:gap-literature}

Three specific gaps motivate what follows.

First, the comparative post-communist literature has largely left Central Asia out of its formal analysis. \textcite{billingsley2010} includes all five republics in her descriptive treatment but restricts the regression sample, excluding the majority of the Central Asian republics; of the five, only Kazakhstan enters the statistical analysis \parencite{spoorenberg2015}. The result is that the region most likely to falsify a general account of post-communist fertility is the region least represented in the tests of it.

Second, the literature stops too early. Billingsley's window closes in 2003, Spoorenberg's principal analyses run to 2009--2011, and the ethnic-differential studies rest on surveys from the 2000s. The turning point documented in Chapter~\ref{sec:trends}, a reversal in dispersion around 2016 and a widening difference thereafter, falls entirely outside every one of these windows. The Uzbek fertility surge after 2017, the largest single movement in the panel assembled here, postdates almost all of the comparative work on the region; the one substantial study of it is a national survey commissioned precisely because existing theory did not anticipate it \parencite{unfpa2023}.

Third, the two literatures use different units and have not been joined. The Central Asian work is micro-level and country-specific; the post-communist comparative work is macro-level and mostly excludes Central Asia. Where a macro treatment of Central Asian fertility exists, it is confined to the region and uses a single economic covariate: \textcite{spoorenberg2015} relates period fertility to GDP per capita in a five-country panel and identifies the reliance on one variable as a limitation. No study places all fourteen comparable post-Soviet states in a single frame, over a window reaching to 2023, with a set of economic controls and inference appropriate to a sample of this size.

This thesis does that, and the ambition is deliberately bounded. It does not attempt to identify why Central Asian fertility is higher; Section~\ref{sec:interpretation} explains why a fourteen-country panel cannot. What it can establish is how large the difference is, whether it has been closing, how much of it moves together with the macroeconomic conditions the comparative literature relies on, and how much is left over for the mechanisms the micro literature describes. Setting that boundary precisely is the contribution, and the size of what remains outside it is the result.

%  No tables or figures.
% =====================================================================

\section{Background and Literature}
\label{sec:background}

\subsection{A common institutional inheritance, and an uncommon starting point}
\label{sec:soviet-inheritance}

Until December 1991 the fourteen countries studied here were administrative units of a single state. They shared a health system, a school curriculum, a language of government, a pension architecture and a family-policy regime, including maternity leave, child benefits and the near-free availability of abortion that made the Soviet Union unusual among industrial societies. If institutional convergence produced demographic convergence, these are the countries where one would expect to see it.

They did not start from a common position, though, and the reason matters for what follows. Fertility in Central Asia through the 1920s and 1930s fluctuated around four to five children per woman, on a par with much of the pre-transitional world. What happened next runs against the usual direction of demographic change. From the early 1940s fertility in the region \emph{rose}, continuing to climb until the early 1960s and holding on a plateau until the early 1970s, before beginning a long decline that bottomed out around 2000 \parencite{spoorenberg2017}.

This happened during exactly the decades when Soviet policy was extending female education, drawing women into the labour force and urbanising the population, the conditions under which the standard account expects fertility to fall. \textcite{spoorenberg2017} rules out the most obvious explanation, showing that marriage was being postponed rather than accelerated during the increase: in Kazakhstan the share of women married at ages 20--24 fell from 78 per cent in 1939 to 62 per cent in 1959. What appears to have driven the rise instead is the removal of biological and behavioural checks on childbearing. Life expectancy in the region rose from around 30--35 years in the mid-1920s to 50--55 by the early 1950s; sterility and infant loss fell with the extension of the health system; and urban migration shortened breastfeeding and therefore birth intervals. Development raised fertility by unleashing childbearing capacity before it began to suppress it, a pattern documented in other early-transition settings \parencite{nag1980,romaniuk1980}.

Two things follow. Central Asia's high fertility at the end of the Soviet period was not simply a pre-modern residue waiting to be modernised away; it was itself partly a product of twentieth-century development. And the region entered independence at a very different point in its fertility history from the European republics, which had been below or near replacement for decades. Whatever else 1991 did, it did not apply a common shock to a common starting position.

\subsection{The transition shock and how it has been explained}
\label{sec:transition}

Fertility fell everywhere in the region after 1991, and in most places it fell hard. Between 1991 and 2000 the total fertility rate declined by 38.6 per cent in Uzbekistan, 33.3 per cent in Kyrgyzstan and 32 per cent in Kazakhstan \parencite{spoorenberg2013}; the European republics fell further and from lower levels, with fifteen post-communist countries recording a rate below 1.3 at least once by 2004 \parencite{billingsley2010}. Populations shrank accordingly. Between 1989 and 2008 Georgia lost nearly 19 per cent of its population and Moldova nearly 18, while Latvia and Estonia lost around 15 per cent each. Over the same period Tajikistan grew by 42 per cent and Uzbekistan by 34 \parencite{brainerd2010}. The demographic divergence this thesis measures in fertility rates was already visible in population totals a decade before the study window opens.

\paragraph{The economic account.} The standard economic framework treats children as time-intensive, so that the shadow price of a child rises with the mother's wage \parencite{becker1981}. On this reading the post-communist decline reflects a rising opportunity cost of childbearing: women's relative wages improved in much of Eastern Europe, and returns to education rose across nearly all transition countries, drawing women into tertiary enrolment and delaying family formation \parencite{brainerd2010}. The account is incomplete on its own terms. \textcite{brainerd2010} notes that women's wages did not improve in Russia or Ukraine, where fertility nonetheless collapsed, and the framework says nothing about uncertainty, which is what the transition mostly supplied. \textcite{kohler2002} argue that the Russian decline of the early and mid-1990s was driven by labour-market crisis and the option value of postponing an irreversible commitment under uncertain conditions. Delay is cheap when the future is unreadable.

\paragraph{Three competing explanations.} \textcite{billingsley2010} organises the debate into three accounts that differ in economic context, in the fertility behaviour they predict, and in what motivates it. The second demographic transition attributes decline to ideational change and self-actualisation, and expects it under economic stability. The postponement transition attributes it to delay of first births as a rational response to uncertainty. The economic-crisis account predicts something different in kind: not delay but stopping, the abandonment of higher-order births under material pressure.

The distinction matters because the two behaviours leave different traces. Postponement depresses the period fertility rate temporarily and may be recovered later; stopping reduces completed family size. Billingsley's central finding is that no single account fits the whole region: postponement explains the bulk of the decline in no more than five countries, all of them in Central Europe, while elsewhere much of the fall happened before postponement began. She also finds the two processes tied to economic conditions in opposite directions, improving conditions with postponement and deteriorating conditions with stopping.

One of her observations is worth carrying forward. Postponement in the post-communist region began from a much younger base than in Southern Europe: mothers' average age at first birth in 1999 was 23.9 years against 28.3 in the Southern European countries where the postponement literature was developed. Delay starting at 24 need not translate into births forgone, which is a caution against reading period declines in this region as permanent.

\paragraph{Policy.} Governments across the region responded with pronatalist measures, of which Russia's 2006 maternity-capital programme is the largest and most studied \parencite{brainerd2010}. \textcite{zakharov2024} assesses it over a long horizon and finds the effect narrower than its reputation: Russian fertility began rising in 2000, six years before the decree, and what the policy most visibly changed was the interval between births rather than the number of them. Period rates rose while births were brought forward, then fell back to 1.50 by 2019, essentially the pre-policy level. He reads this as the same inflation-and-correction pattern produced by the pronatalist campaign of the 1980s, and concludes that two decades of policy did not halt the underlying modernisation of Russian fertility. The wider evidence points the same way: financial incentives tend to move the timing of childbearing more than its quantum \parencite{gauthier2007}.

\subsection{What is known about Central Asia}
\label{sec:ca-literature}

The Central Asian literature is smaller, more recent, and built almost entirely on country cases rather than regional comparisons.

The starting point is that the fertility recovery after 2000 is real. \textcite{spoorenberg2015} tests the possibility that it reflects improving completeness of birth registration by cross-comparing vital-registration rates against census and survey estimates computed by independent methods, and rules the artefact explanation out, with a reservation for Turkmenistan, where independent data are not accessible. This matters more than it might appear: in a region where one country has not held a census since 1989, establishing that a demographic movement is not a recording artefact is a precondition for studying it at all.

Beyond that, four strands of explanation have developed, and because Chapter~\ref{sec:discussion} works through them in detail, they are only mapped here.

The first is population composition. Fertility differs by about one child per woman between women of the titular ethnicities and women of European origin, consistently across countries and data sources, and post-1991 emigration changed the ethnic composition of these populations substantially \parencite{spoorenberg2013,agadjanian2013}. Holding ethnic-specific fertility constant and letting composition alone change reproduces a large share of the observed national increases \parencite{spoorenberg2015}. The second is marriage: near-universal but with rising age at first union, and sharply differentiated by education \parencite{dommaraju2008}. The third is culture, understood not as individual belief but as group-level norms and expectations. \textcite{agadjanian2022} distinguish the normative propensities a group transmits from its position in a society's hierarchy, and find the two predicting different things: completed fertility follows the first, desires for further children the second. \textcite{kumo2024} approach the same territory from the religion side, modelling a chain from religiosity through gender-role beliefs to family size and asking how far Soviet institutions interrupted it. The fourth is the economic context, which \textcite{spoorenberg2015} finds positively associated with fertility in a five-country panel, alongside a shifting tempo effect accounting for part of the rise in Kyrgyzstan and rather less of it in Kazakhstan. Labour migration belongs here too: remittances reach a fifth or more of GDP in Tajikistan and Kyrgyzstan, and their relationship with fertility across the post-communist world is negative on average but heterogeneous \parencite{tokhirov2024}.

Two features of this literature are worth naming. It is predominantly micro-level, working with survey and census microdata, which is what allows it to identify mechanisms a country-level design cannot. And it is largely internal to the region, comparing Central Asian countries with one another or ethnic groups within them rather than against the rest of the post-Soviet space. \textcite{nedoluzhko2012} is the most explicitly comparative, and finds country of residence mattering more for fertility than ethnicity, even among co-ethnics living either side of a Central Asian border.

\subsection{The gap}
\label{sec:gap-literature}

Three specific gaps motivate what follows.

First, the comparative post-communist literature has largely left Central Asia out of its formal analysis. \textcite{billingsley2010} includes all five republics in her descriptive treatment but restricts the regression sample, excluding the majority of the Central Asian republics; of the five, only Kazakhstan enters the statistical analysis \parencite{spoorenberg2015}. The result is that the region most likely to falsify a general account of post-communist fertility is the region least represented in the tests of it.

Second, the literature stops too early. Billingsley's window closes in 2003, Spoorenberg's principal analyses run to 2009--2011, and the ethnic-differential studies rest on surveys from the 2000s. The turning point documented in Chapter~\ref{sec:trends}, a reversal in dispersion around 2016 and a widening difference thereafter, falls entirely outside every one of these windows. The Uzbek fertility surge after 2017, the largest single movement in the panel assembled here, postdates almost all of the comparative work on the region; the one substantial study of it is a national survey commissioned precisely because existing theory did not anticipate it \parencite{unfpa2023}.

Third, the two literatures use different units and have not been joined. The Central Asian work is micro-level and country-specific; the post-communist comparative work is macro-level and mostly excludes Central Asia. Where a macro treatment of Central Asian fertility exists, it is confined to the region and uses a single economic covariate: \textcite{spoorenberg2015} relates period fertility to GDP per capita in a five-country panel and identifies the reliance on one variable as a limitation. No study places all fourteen comparable post-Soviet states in a single frame, over a window reaching to 2023, with a set of economic controls and inference appropriate to a sample of this size.

This thesis does that, and the ambition is deliberately bounded. It does not attempt to identify why Central Asian fertility is higher; Section~\ref{sec:interpretation} explains why a fourteen-country panel cannot. What it can establish is how large the difference is, whether it has been closing, how much of it moves together with the macroeconomic conditions the comparative literature relies on, and how much is left over for the mechanisms the micro literature describes. Setting that boundary precisely is the contribution, and the size of what remains outside it is the result.

%  No tables or figures.
% =====================================================================

\section{Background and Literature}
\label{sec:background}

\subsection{A common institutional inheritance, and an uncommon starting point}
\label{sec:soviet-inheritance}

Until December 1991 the fourteen countries studied here were administrative units of a single state. They shared a health system, a school curriculum, a language of government, a pension architecture and a family-policy regime, including maternity leave, child benefits and the near-free availability of abortion that made the Soviet Union unusual among industrial societies. If institutional convergence produced demographic convergence, these are the countries where one would expect to see it.

They did not start from a common position, though, and the reason matters for what follows. Fertility in Central Asia through the 1920s and 1930s fluctuated around four to five children per woman, on a par with much of the pre-transitional world. What happened next runs against the usual direction of demographic change. From the early 1940s fertility in the region \emph{rose}, continuing to climb until the early 1960s and holding on a plateau until the early 1970s, before beginning a long decline that bottomed out around 2000 \parencite{spoorenberg2017}.

This happened during exactly the decades when Soviet policy was extending female education, drawing women into the labour force and urbanising the population, the conditions under which the standard account expects fertility to fall. \textcite{spoorenberg2017} rules out the most obvious explanation, showing that marriage was being postponed rather than accelerated during the increase: in Kazakhstan the share of women married at ages 20--24 fell from 78 per cent in 1939 to 62 per cent in 1959. What appears to have driven the rise instead is the removal of biological and behavioural checks on childbearing. Life expectancy in the region rose from around 30--35 years in the mid-1920s to 50--55 by the early 1950s; sterility and infant loss fell with the extension of the health system; and urban migration shortened breastfeeding and therefore birth intervals. Development raised fertility by unleashing childbearing capacity before it began to suppress it, a pattern documented in other early-transition settings \parencite{nag1980,romaniuk1980}.

Two things follow. Central Asia's high fertility at the end of the Soviet period was not simply a pre-modern residue waiting to be modernised away; it was itself partly a product of twentieth-century development. And the region entered independence at a very different point in its fertility history from the European republics, which had been below or near replacement for decades. Whatever else 1991 did, it did not apply a common shock to a common starting position.

\subsection{The transition shock and how it has been explained}
\label{sec:transition}

Fertility fell everywhere in the region after 1991, and in most places it fell hard. Between 1991 and 2000 the total fertility rate declined by 38.6 per cent in Uzbekistan, 33.3 per cent in Kyrgyzstan and 32 per cent in Kazakhstan \parencite{spoorenberg2013}; the European republics fell further and from lower levels, with fifteen post-communist countries recording a rate below 1.3 at least once by 2004 \parencite{billingsley2010}. Populations shrank accordingly. Between 1989 and 2008 Georgia lost nearly 19 per cent of its population and Moldova nearly 18, while Latvia and Estonia lost around 15 per cent each. Over the same period Tajikistan grew by 42 per cent and Uzbekistan by 34 \parencite{brainerd2010}. The demographic divergence this thesis measures in fertility rates was already visible in population totals a decade before the study window opens.

\paragraph{The economic account.} The standard economic framework treats children as time-intensive, so that the shadow price of a child rises with the mother's wage \parencite{becker1981}. On this reading the post-communist decline reflects a rising opportunity cost of childbearing: women's relative wages improved in much of Eastern Europe, and returns to education rose across nearly all transition countries, drawing women into tertiary enrolment and delaying family formation \parencite{brainerd2010}. The account is incomplete on its own terms. \textcite{brainerd2010} notes that women's wages did not improve in Russia or Ukraine, where fertility nonetheless collapsed, and the framework says nothing about uncertainty, which is what the transition mostly supplied. \textcite{kohler2002} argue that the Russian decline of the early and mid-1990s was driven by labour-market crisis and the option value of postponing an irreversible commitment under uncertain conditions. Delay is cheap when the future is unreadable.

\paragraph{Three competing explanations.} \textcite{billingsley2010} organises the debate into three accounts that differ in economic context, in the fertility behaviour they predict, and in what motivates it. The second demographic transition attributes decline to ideational change and self-actualisation, and expects it under economic stability. The postponement transition attributes it to delay of first births as a rational response to uncertainty. The economic-crisis account predicts something different in kind: not delay but stopping, the abandonment of higher-order births under material pressure.

The distinction matters because the two behaviours leave different traces. Postponement depresses the period fertility rate temporarily and may be recovered later; stopping reduces completed family size. Billingsley's central finding is that no single account fits the whole region: postponement explains the bulk of the decline in no more than five countries, all of them in Central Europe, while elsewhere much of the fall happened before postponement began. She also finds the two processes tied to economic conditions in opposite directions, improving conditions with postponement and deteriorating conditions with stopping.

One of her observations is worth carrying forward. Postponement in the post-communist region began from a much younger base than in Southern Europe: mothers' average age at first birth in 1999 was 23.9 years against 28.3 in the Southern European countries where the postponement literature was developed. Delay starting at 24 need not translate into births forgone, which is a caution against reading period declines in this region as permanent.

\paragraph{Policy.} Governments across the region responded with pronatalist measures, of which Russia's 2006 maternity-capital programme is the largest and most studied \parencite{brainerd2010}. \textcite{zakharov2024} assesses it over a long horizon and finds the effect narrower than its reputation: Russian fertility began rising in 2000, six years before the decree, and what the policy most visibly changed was the interval between births rather than the number of them. Period rates rose while births were brought forward, then fell back to 1.50 by 2019, essentially the pre-policy level. He reads this as the same inflation-and-correction pattern produced by the pronatalist campaign of the 1980s, and concludes that two decades of policy did not halt the underlying modernisation of Russian fertility. The wider evidence points the same way: financial incentives tend to move the timing of childbearing more than its quantum \parencite{gauthier2007}.

\subsection{What is known about Central Asia}
\label{sec:ca-literature}

The Central Asian literature is smaller, more recent, and built almost entirely on country cases rather than regional comparisons.

The starting point is that the fertility recovery after 2000 is real. \textcite{spoorenberg2015} tests the possibility that it reflects improving completeness of birth registration by cross-comparing vital-registration rates against census and survey estimates computed by independent methods, and rules the artefact explanation out, with a reservation for Turkmenistan, where independent data are not accessible. This matters more than it might appear: in a region where one country has not held a census since 1989, establishing that a demographic movement is not a recording artefact is a precondition for studying it at all.

Beyond that, four strands of explanation have developed, and because Chapter~\ref{sec:discussion} works through them in detail, they are only mapped here.

The first is population composition. Fertility differs by about one child per woman between women of the titular ethnicities and women of European origin, consistently across countries and data sources, and post-1991 emigration changed the ethnic composition of these populations substantially \parencite{spoorenberg2013,agadjanian2013}. Holding ethnic-specific fertility constant and letting composition alone change reproduces a large share of the observed national increases \parencite{spoorenberg2015}. The second is marriage: near-universal but with rising age at first union, and sharply differentiated by education \parencite{dommaraju2008}. The third is culture, understood not as individual belief but as group-level norms and expectations. \textcite{agadjanian2022} distinguish the normative propensities a group transmits from its position in a society's hierarchy, and find the two predicting different things: completed fertility follows the first, desires for further children the second. \textcite{kumo2024} approach the same territory from the religion side, modelling a chain from religiosity through gender-role beliefs to family size and asking how far Soviet institutions interrupted it. The fourth is the economic context, which \textcite{spoorenberg2015} finds positively associated with fertility in a five-country panel, alongside a shifting tempo effect accounting for part of the rise in Kyrgyzstan and rather less of it in Kazakhstan. Labour migration belongs here too: remittances reach a fifth or more of GDP in Tajikistan and Kyrgyzstan, and their relationship with fertility across the post-communist world is negative on average but heterogeneous \parencite{tokhirov2024}.

Two features of this literature are worth naming. It is predominantly micro-level, working with survey and census microdata, which is what allows it to identify mechanisms a country-level design cannot. And it is largely internal to the region, comparing Central Asian countries with one another or ethnic groups within them rather than against the rest of the post-Soviet space. \textcite{nedoluzhko2012} is the most explicitly comparative, and finds country of residence mattering more for fertility than ethnicity, even among co-ethnics living either side of a Central Asian border.

\subsection{The gap}
\label{sec:gap-literature}

Three specific gaps motivate what follows.

First, the comparative post-communist literature has largely left Central Asia out of its formal analysis. \textcite{billingsley2010} includes all five republics in her descriptive treatment but restricts the regression sample, excluding the majority of the Central Asian republics; of the five, only Kazakhstan enters the statistical analysis \parencite{spoorenberg2015}. The result is that the region most likely to falsify a general account of post-communist fertility is the region least represented in the tests of it.

Second, the literature stops too early. Billingsley's window closes in 2003, Spoorenberg's principal analyses run to 2009--2011, and the ethnic-differential studies rest on surveys from the 2000s. The turning point documented in Chapter~\ref{sec:trends}, a reversal in dispersion around 2016 and a widening difference thereafter, falls entirely outside every one of these windows. The Uzbek fertility surge after 2017, the largest single movement in the panel assembled here, postdates almost all of the comparative work on the region; the one substantial study of it is a national survey commissioned precisely because existing theory did not anticipate it \parencite{unfpa2023}.

Third, the two literatures use different units and have not been joined. The Central Asian work is micro-level and country-specific; the post-communist comparative work is macro-level and mostly excludes Central Asia. Where a macro treatment of Central Asian fertility exists, it is confined to the region and uses a single economic covariate: \textcite{spoorenberg2015} relates period fertility to GDP per capita in a five-country panel and identifies the reliance on one variable as a limitation. No study places all fourteen comparable post-Soviet states in a single frame, over a window reaching to 2023, with a set of economic controls and inference appropriate to a sample of this size.

This thesis does that, and the ambition is deliberately bounded. It does not attempt to identify why Central Asian fertility is higher; Section~\ref{sec:interpretation} explains why a fourteen-country panel cannot. What it can establish is how large the difference is, whether it has been closing, how much of it moves together with the macroeconomic conditions the comparative literature relies on, and how much is left over for the mechanisms the micro literature describes. Setting that boundary precisely is the contribution, and the size of what remains outside it is the result.
